\Cref{fig:1-challenges} connects three challenges to the mechanisms and
evidence used to address them.

First, the primary codebook has useful mathematical structure that a direct
GPU port can overlook. Each eight-bit subspace code indexes 256 distances.
Its Cartesian coordinates permit an exact 16+16 factorisation, but more
aggressive splitting adds lookups and arithmetic. The challenge is to exploit
separability without assuming that the smallest table gives the fastest
scan. Candidate-parallel word access and residual lane cooperation also need
different data mappings.

Second, a query's primary winners can be concentrated in any part of its
routed pool. Parallel groups cannot safely retain only a fraction of the
required global survivors. Exact local retention, a global merge and a
shared complete-score bound provide a principled route, but add buffers,
staging and synchronisation. A large batch can already supply enough blocks;
an imbalanced workload can still benefit from more groups. The challenge
couples correctness with a workload-dependent execution decision. Selecting
IVF lists remains a separate primitive from retaining database candidates.

Third, a fixed survivor multiplier can repeatedly refine irrelevant vectors,
while a small multiplier can remove a true neighbour too early. Agreement
with a large-budget reference is available without GT, but agreement does
not repair the reference's own errors or guarantee future-query recall.
Choosing a sample and budget also consumes time. The challenge is to evaluate
selection quality, calibration cost and reuse together.

\par\medskip\noindent\begin{minipage}{\linewidth}
\centering\captionsetup{hypcap=false}
\resizebox{\linewidth}{!}{\begin{tikzpicture}[
 font=\small\linespread{1}\selectfont,
 box/.style={draw=ThesisNavy!65,rounded corners=2pt,align=center,
 text width=3.5cm,minimum height=1.65cm,inner sep=4pt},
 arr/.style={-{Latex[length=2mm]},draw=ThesisNavy!80,thick}]
\node[font=\bfseries] at (0,1.25) {Challenge};
\node[font=\bfseries] at (4.6,1.25) {Mechanism};
\node[font=\bfseries] at (9.2,1.25) {Evidence};
\node[box] (a) at (0,0) {Primary codes\\Broad scan cost};
\node[box,fill=ThesisNavy!7] (b) at (4.6,0) {Cartesian tables\\Packed word scan};
\node[box] (c) at (9.2,0) {Table / load tests\\Score checks};
\node[box] (d) at (0,-2.2) {Uneven query work\\Global top-$c$};
\node[box,fill=ThesisNavy!7] (e) at (4.6,-2.2) {Grouped scan\\Shared bound\\Global merge};
\node[box] (f) at (9.2,-2.2) {Output checks\\Batch / skew tests};
\node[box] (g) at (0,-4.4) {Refinement effort\\Quality / reuse cost};
\node[box,fill=ThesisNavy!7] (h) at (4.6,-4.4) {Sampled agreement\\Workload budget};
\node[box] (i) at (9.2,-4.4) {Held-out quality\\Calibration cost};
\draw[arr] (a)--(b);\draw[arr] (b)--(c);
\draw[arr] (d)--(e);\draw[arr] (e)--(f);
\draw[arr] (g)--(h);\draw[arr] (h)--(i);
\end{tikzpicture}
}
\captionof{figure}[Problem, mechanism and evidence]{Problem, mechanism and evidence. The arrows organise
research questions; they do not imply that independently measured speedups
multiply or that an unlabelled agreement rule guarantees GT recall.}
\label{fig:1-challenges}
\end{minipage}\par\medskip
